\section{Preliminary Results and Implementation Status}
The Ecogent system proposed in this research offers a complete framework for AI automation that prioritizes local processing over the cloud. To prove that these ideas work, an initial prototype has been successfully built. The prototype's source code is publicly available and actively maintained at \url{https://github.com/samratpro/Ecogent}. This early version demonstrates the core routing system, the local memory, and specific web browser workflows. While the complete system (including all specialized agents mentioned in the methodology) will be built during the main research phase, this prototype proves that the local-first approach works and saves significant money.\\

\subsection*{Current Prototype Capabilities}
At this stage of development, the Ecogent prototype has reached its first major goals. It now has a fully working, local-first management pipeline.\\

\textbf{Completed Tasks:}
\begin{itemize}
    \item \textbf{Local Supervisor:} Created a free, fast tool that perfectly sorts user requests based on how difficult they are.
    \item \textbf{Local Memory (ChromaDB):} Set up a local database to quickly search for tools, save browser actions, and remember past sessions.
    \item \textbf{Browser Record and Replay:} Built a web agent that records actions and replays them perfectly, skipping expensive AI costs for repeated tasks.
    \item \textbf{Automatic Self-Fixing:} Designed a recovery tool that instantly fixes broken browser actions if a website changes.
    \item \textbf{Small Local AI:} Added a tiny AI model that can answer simple questions without using any expensive cloud tokens.
    \item \textbf{Token-Saving Execution:} Replaced heavy cloud memory systems with strict token limits and smart error shortening.
\end{itemize}

\textbf{Pending Tasks (Future Work Roadmap):}
\begin{itemize}
    \item \textbf{Dynamic Cloud Routing:} Building a smarter local filter to pick the cheapest and best cloud model (like Llama-3 or GPT-4) based on exactly what is needed.
    \item \textbf{Coding Agent:} Adding a dedicated software engineering agent that uses a smart code map.
    \item \textbf{Testing Agent:} Building the token-optimized QA testing agent (currently partially implemented).
    \item \textbf{OS Agent:} Creating the continuous background terminal system for managing folders and commands (currently partially implemented).
    \item \textbf{Data and API Agents:} Developing specialized agents to handle large spreadsheets (CSV/Excel) and web connections.
\end{itemize}

\textbf{Project Timeline (Pre-Defense Schedule):}
To ensure the remaining components are systematically developed and tested, the following timeline is proposed:\\
\begin{itemize}
    \item \textbf{Phase 1 (Completed):} Prototype the core multi-tier routing system, NLP supervisor, and local memory (ChromaDB).
    \item \textbf{Phase 2 (Upcoming):} Fully develop the Testing Agent and OS Agent to support token truncation and background processes.
    \item \textbf{Phase 3 (Upcoming):} Implement the Coding Agent and its Semantic Code Map layer.
    \item \textbf{Phase 4 (Final Evaluation):} Run rigorous benchmarking using SWE-bench and WebArena to finalize token and cost reduction metrics.
\end{itemize}

\section{Expected Outcomes}
The expected results of this research are:\\
\begin{itemize}
    \item \textbf{No More Wasted Tokens:} Stops the system from reading useless, long histories by only sending the most important, critical information to the AI.
    \item \textbf{Lower Cost and Faster Speed:} Greatly reduces API costs and waiting times by handling simple questions and repeated tasks directly on the local computer.
    \item \textbf{Reliable Performance:} Improves the speed and trust of the system by replacing messy conversations with clear, step-by-step JSON plans.
    \item \textbf{Smart Memory Reuse:} Reuses old solutions and browser recordings from the local database, stopping the AI from wasting time thinking about a problem it has already solved.
    \item \textbf{Better Hardware Management:} Gives the user more control over how much power the agent uses, keeping local processing lightweight and smooth.
\end{itemize}
